\section*{Erd\H{o}s Problem 403}

\paragraph{Statement and result.}
Does \(2^m=a_1!+\cdots+a_k!\), with distinct positive indices
\(a_1<\cdots<a_k\), have only finitely many solutions?
Yes. We reconstruct the known Frankl--Lin result using an explicit finite
integer computation: the possible exponents are
\[
 m=0,1,3,5,7.
 \tag{1}
\]
The largest solution is \(128=2!+3!+5!\).
The small computation below is part of the proof; no Lean verification
of this particular reconstruction is claimed.

\paragraph{Reduction to a bounded computation.}
For \(m\geq1\), parity excludes \(1!\). Reducing modulo \(3\)
then forces \(2!\) to occur, since every factorial with index at least
three is divisible by three. We prove that no sum of distinct factorials
containing \(2!\) and no \(1!\) is divisible by \(2^{255}\).

Put \(S_2=\{2\}\). For \(3\leq n\leq255\), define
\[
 e_n=\min(255,v_2((n+1)!)),\qquad
 S_n=\{x:x=s\text{ or }s+n!,\ s\in S_{n-1},\ 2^{e_n}\mid x\}.
 \tag{2}
\]
If a full factorial sum is divisible by \(2^{255}\), its partial
sum through \(n!\) must lie in \(S_n\): all omitted future factorials
are divisible by \(2^{e_n}\). This follows inductively, including
zero choices at missing indices. At \(n=255\), all later factorials
are multiples of \(2^{255}\). Therefore \(S_{255}=\varnothing\)
excludes every full sum, including those with indices exceeding \(255\).

\paragraph{Exact certificate.}
The following complete Python code uses only unbounded integer arithmetic.
The identity \(v_2(t!)=t-\operatorname{popcount}(t)\) evaluates (2).
\begin{verbatim}
from math import factorial
S, F = {2}, 2
peak = 1
for n in range(3, 256):
    F *= n
    e = min(255, n + 1 - (n + 1).bit_count())
    S = {x for s in S for x in (s, s + F)
         if x % (1 << e) == 0}
    peak = max(peak, len(S))
assert peak == 496
assert S == set()

facts = [factorial(i) for i in range(1, 256)]
answers = []
for m in range(255):
    remainder, indices = 1 << m, []
    for i in range(255, 0, -1):
        if facts[i - 1] <= remainder:
            remainder -= facts[i - 1]
            indices.append(i)
    if remainder == 0:
        answers.append((m, sorted(indices)))
assert answers == [(0,[1]), (1,[2]), (3,[2,3]),
                   (5,[2,3,4]), (7,[2,3,5])]
\end{verbatim}
For additional audit points, the cardinalities \(|S_n|\) for
\(n=250,251,252,253,254,255\) are respectively
\(104,57,114,114,228,0\). The code was run successfully.

The first assertion proves \(m\leq254\). Also \(256!>2^{254}\),
so indices at least \(256\) cannot occur in that remaining range.
The greedy check is exhaustive because
\(\sum_{i=1}^{t-1}i!<t!\) for every \(t\geq2\).
Thus the largest available factorial must be selected if it does not
exceed the remainder; otherwise no combination of smaller ones can
replace it. This proves both the enumeration and (1).

\paragraph{Index convention and attribution.}
The argument uses positive indices. If \(0!\) is also allowed, finiteness
and the largest solution remain the same. For \(m\geq3\), parity says
that \(0!\) and \(1!\) occur together or neither does. The contribution
from indices \(0,1,2\) is then \(0,2,\) or \(4\). Modulo \(3\) rules out
zero. A contribution of \(4\), together with the optional \(3!=6\),
is \(4\) or \(2\pmod8\), so is impossible. The contribution is therefore
exactly \(2\), and the same computation applies. The finitely many
smaller exponents cause no issue.

The known result is attributed to Frankl and Lin at
\url{https://www.erdosproblems.com/403}, checked 24 September 2026.
This is a reproducible reconstruction of that result, not a new discovery.

\paragraph{Completion estimate.}
The infinite search is rigorously reduced to, and settled by, the displayed finite certificate.
\textbf{COMPLETION: 100\%}
